\documentclass[noanswers, nolegalese, 11pt]{examjh}
% \documentclass[answers, nolegalese, 11pt]{examjh}
\usepackage{../../../../computingfonts}
\usepackage{../../../../contentmacros}
\usepackage{../../../../grammar}
\usepackage{graphicx}

\setlength{\parindent}{0em}\setlength{\parskip}{0.5ex}
\pagestyle{empty}
\begin{document}\thispagestyle{empty}
\makebox[\textwidth]{Questions for Five.V\hfill  From \textit{Theory of Computation}, by Hef{}feron}\vspace{-1ex}
\makebox[\textwidth]{\hbox{}\hrulefill\hbox{}}



\begin{questions}
\question
There are $n$ radio stations and we must assign each a radio frequency.
There are $f$~many frequencies available.  
To avoid interference, stations closer than $100$~km cannot use the
same frequency, so we have a list of distances $D(i,j)$
and want to know if there is a set of assignments that meets the restricitions.
Restate this as a satisfiability problem.
\begin{parts}
\part
Give a propositional logic statement for the case of $n=3$ and $f=2$. 
\begin{solution}
The $n=3$ and $f=2$ case shows what to do.
The condition $D(0,1)\geq 100$ is either True or False; call it $P$.
Similarly, let $Q$ be the condition $D(0,2)\geq 100$,
and let $R$ be the condition $D(1,2)\geq 100$.
Then this statement is satisfiable if and only if there is a frequency 
assignment that meets the restrictions.
\begin{equation*}
  (P\wedge Q)\vee(P\wedge R)\vee(Q\wedge R)
\end{equation*}
\end{solution}

\part For the propositional logic statement in the $n,f$ general case, 
how many $\wedge$-collected clauses
are there?
\begin{solution}
$\binom{n}{f}$
\end{solution}
\end{parts}

\question
A judge must form a jury of twelve people. 
The law requires that no one on the jury know any other member. 
There are $n$ candidates and the judge has a list of Booleans
$\operatorname{knows}(i,j)$.
Restate this as a satisfiability problem.
\begin{parts}
\part
Give a propositional logic statement for the case of 
forming juries of three people chosen from a pool of five.
\begin{solution}
For the hint, the judge has clauses that look like this.
\begin{align*}
  &\bigl[(\neg\operatorname{knows}(0,1)\wedge\neg\operatorname{knows}(0,2)\wedge\neg\operatorname{knows}(1,2)\bigr]  \\
  &\quad\vee\,
  \bigl[(\neg\operatorname{knows}(0,1)\wedge\neg\operatorname{knows}(0,3)\wedge\neg\operatorname{knows}(1,3)\bigr]  \\
  &\quad\vee\qquad\smash{\vdots{}}           \\
  &\quad\vee\,\bigl[(\neg\operatorname{knows}(3,4)\wedge\neg\operatorname{knows}(3,5)\wedge\neg\operatorname{knows}(4,5)\bigr]
\end{align*}
\end{solution}
\part How many $\wedge$-connected clauses are there?
\begin{solution}
There is a clause for each triple of $1$'s
\begin{center}
\begin{tabular}{*{5}{c}}
  \textit{candidate $0$}
    &\textit{candidate $1$}
    &\textit{candidate $2$}
    &\textit{candidate $3$}
    &\textit{candidate $4$}  \\
  \hline
  $1$ &$1$ &$1$ &$0$ &$0$ \\
  $1$ &$1$ &$0$ &$1$ &$0$ \\
  $1$ &$1$ &$0$ &$0$ &$1$ \\
  $1$ &$0$ &$1$ &$1$ &$0$ \\
  $1$ &$0$ &$1$ &$0$ &$1$ \\
  $1$ &$0$ &$0$ &$1$ &$1$ \\
  $0$ &$1$ &$1$ &$1$ &$0$ \\
  $0$ &$1$ &$1$ &$0$ &$1$ \\
  $0$ &$1$ &$0$ &$1$ &$1$ \\
  $0$ &$0$ &$1$ &$1$ &$1$ \\
\end{tabular}
\end{center}
There are $\binom{5}{3}=10$ many clauses.
\end{solution}

\part In the general $n$-candidate case, how many clauses are there?
\begin{solution}
There are $\binom{n}{12}$~many clauses.
\end{solution}
\end{parts}

\question
Recall that we say a graph 
is \textit{$k$-colored} if the vertices are
broken into $k$-many disjoint sets so that no two vertices in a 
set are connected
(these sets are called \textit{colors} because of the problem's history).
The \textit{$k$-colorability problem} is, given a graph, to determine if it
can be $k$-colored.
Restate this as a satisfiability problem.
\begin{parts}
\part
Consider a graph $\mathcal{G}$ with five vertices, 
$\mathcal{N}=\set{v_0,\ldots v_4}$.
How many ways are there to break the vertices of $\mathcal{N}$ into 
two colors?
(This pays no attention to whether they are connected.)
\begin{solution}
Take a subset $C_0\subseteq\mathcal{N}$, and then the
the complement is the other color, $C_1=\mathcal{N}-C_0$.
Because $\sizeof{\mathcal{N}}=5$, there are $2^5=32$ many possible such subsets
(although there is a redundancy such as with 
$C_0=\set{v_0,v_1},C_1=\set{v_2,v_3,v_4}$ and
$C_0=\set{v_2,v_3,v_4},C_1=\set{v_0,v_1}$).
\end{solution}
\part Assume that you have a Boolean function $e(i,j)$ that is True if and only
if $v_i$ is connected to~$v_j$.
Give a propositional logic statement that is satisfiable if
and only if this graph is $2$-colorable.
You can simply produce the clause for the 
$C_0=\set{v_0,v_1}$ and $C_1=\set{v_2,v_3,v_4}$ potential coloring.
\begin{solution}
For the 
$C_0=\set{v_0,v_1}$ and $C_1=\set{v_2,v_3,v_4}$ potential coloring
we get this clause.
\begin{equation*}
  \neg e(0,1)
  \,\wedge\, 
  \neg e(2,3) 
  \,\wedge\,
  \neg e(2,4) 
  \,\wedge\,
  \neg e(3,4) 
\end{equation*}
\end{solution}
\part
Outline how to get a suitable propositional logic statement for the
$k$-colorability problem.
\end{parts}

\question
Give an algorithm to solve the frequency assignment problem
suitable for a nondeterministic Turing machine.

\question
Give an algorithm to solve the jury selection problem
suitable for a nondeterministic Turing machine.

\question
Give an algorithm to solve the $k$-colorable problem
suitable for a nondeterministic Turing machine.


\end{questions}
\end{document}
